\bookchapter{第三十章：投屏、演示與車載互聯}

\chapterintro{手機螢幕再大都係7吋，但你嘅電視有55吋、你嘅投影機有100吋。將手機畫面「掟」到大螢幕上，體驗係完全唔同維度嘅。呢一章講清楚投屏、演示同車載互聯呢三個高頻場景，幫你將手機嘅能力延伸到更大嘅螢幕上。}

\booksection{一、無線投屏三大陣營——唔好搞混}

無線投屏睇落都係「手機畫面到電視上」，但底層技術有三種，互唔通用：

\begin{center}
\begin{tabularx}{\linewidth}{|X|X|X|}
\hline
\textbf{技術} & \textbf{邊個用} & \textbf{特點} \tabularnewline
\hline
AirPlay & iPhone/iPad/Mac & 蘋果生態，體驗最流暢。需要Apple TV或支援AirPlay嘅電視 \tabularnewline
Miracast & 安卓手機 & Wi-Fi Direct直連，唔需要路由器。安卓手機「無線投屏」預設用呢個 \tabularnewline
DLNA & 所有平臺 & 淨係推送媒體檔案（影片/音樂/相片），唔鏡像螢幕。適合推影片到電視上播放 \tabularnewline
\hline
\end{tabularx}
\end{center}

\booksubsection{各平臺操作方法：}

\begin{itemize}
  \item \textbf{iPhone}：控制中心 > 螢幕鏡像 > 選擇你嘅電視/Apple TV。如果電視唔支援AirPlay，可以買一個Apple TV盒子或支援AirPlay嘅電視棒。
  \item \textbf{安卓}：下拉通知欄 > 搵到「無線投屏」/「智能投屏」/「多屏互動」按鈕 > 選擇電視。需要手機同電視連同一個Wi-Fi。
  \item \textbf{萬能方案}：如果電視太老唔支援任何無線投屏，買一個Chromecast或小米盒子/華為盒子插到電視HDMI口上，就能接收投屏喇。
\end{itemize}

\booksection{二、有線投屏——更穩定、零延遲}

無線投屏方便但有延遲，打機同做演示時體驗唔好。有線投屏能解決呢個問題：

\begin{itemize}
  \item \textbf{iPhone}：需要一根Lightning/USB-C轉HDMI嘅轉換線（約100-200蚊）。插上後手機畫面直接輸出到電視/投影機嘅HDMI口。
  \item \textbf{安卓}：部分手機支援USB-C轉HDMI直連（需要手機支援DP Alt Mode，大部分旗艦都支援）。買一根USB-C to HDMI線即可。
  \item \textbf{會議室場景}：如果你經常需要用手機做演示，建議喺袋入面便常備一根轉HDMI線。會議室嘅電視/投影機幾乎都有HDMI口，插上就能用，比折騰無線投屏靠譜得多。
\end{itemize}

\booksection{三、手機變電腦——桌面模式}

部分安卓手機支援「桌面模式」——連上顯示器後，手機會輸出一個類似電腦桌面嘅介面，可以用滑鼠鍵盤操作：

\begin{itemize}
  \item \textbf{Samsung DeX}：三星旗艦嘅招牌功能。連接顯示器後彈出桌面介面，可以多視窗操作，體驗最接近真正嘅電腦。配合藍牙鍵鼠，出差可以唔帶筆記本。
  \item \textbf{華為桌面模式}：華為Mate系列/Pura系列連接顯示器後亦有類似功能。支援多視窗、檔案管理、辦公App。
  \item \textbf{摩托羅拉Ready For}：Moto旗艦亦支援，功能類似DeX。
  \item \textbf{適用場景}：輕辦公（回郵件、改文檔）、演示PPT、大屏睇片。唔適合重度Excel操作或專業設計。
\end{itemize}

\booksection{四、車載互聯——揸車用手機嘅正確方式}

揸車時操作手機係違法嘅，但導航、聽音樂、接電話又係剛需。車載互聯系統等你喺車機螢幕上使用手機嘅核心功能，手機可以放埋一邊唔碰：

\begin{itemize}
  \item \textbf{Apple CarPlay}：iPhone用戶嘅首選。透過USB或無線連接車機，車機螢幕上出現iPhone嘅精簡介面——導航（高德/百度）、音樂、電話、資訊、播客。Siri語音控制全程免手。
  \item \textbf{Android Auto}：安卓用戶對應方案。功能與CarPlay類似，支援Google Maps、Spotify等。國內因Google服務限制，部分功能唔可用。
  \item \textbf{華為HiCar}：華為自研方案，國內車型支援率高。除咗導航音樂，仲能用華為視像通話、泊車影像融合等特色功能。
  \item \textbf{百度CarLife}：跨平臺方案，iOS同安卓都能用。功能較少但相容性廣，好多舊車型淨係支援CarLife。
\end{itemize}

\begin{tcbwarn}{安全提醒}
無論邊種車載互聯方案，揸車時都應盡量使用語音控制。導航提前設好目的地，音樂提前選好歌單，唔好喺行駛中操作車機螢幕。

\end{tcbwarn}

\booksection{五、會議室投屏實戰——唔再手忙腳亂}

開會係投屏最高頻嘅場景，亦係最容易出問題嘅場景。以下係幾條實戰經驗：

\begin{itemize}
  \item \textbf{提前到場測試}：永遠唔好等到會議開始先至連投屏。提前10分鐘到場，測試好連接方式。Wi-Fi投屏失敗嘅概率比你諗嘅高得多。
  \item \textbf{準備有線方案做Plan B}：喺袋入面便放一根轉HDMI線。無線投屏連唔上嘅時候，有線就係救命稻草。
  \item \textbf{關閉通知預覽}：投屏時你嘅微信消息會顯示喺會議室大屏上。投屏前喺設定入面便關閉「鎖屏通知內容預覽」，或者開啟「演示模式」（部分手機支援，投屏時自動隱藏通知）。
  \item \textbf{用「演示文檔」而唔係「操作手機」}：如果你要展示PPT，直接將PPT存到手機入面便用WPS打開全屏播放，而唔係投屏整個手機桌面。咁樣更專業、更流暢。
\end{itemize}

\begin{tcbpractice}{動手：完成一次手機到電視嘅投屏}
\begin{enumerate}
  \item 確保手機同電視連同一個Wi-Fi。
  \item 用手機打開一張相片或一個影片。
  \item 嘗試無線投屏（iPhone用螢幕鏡像，安卓用無線投屏按鈕）。
  \item 如果失敗，檢查電視係咪支援對應協議，或嘗試DLNA推送（喺相冊/影片App入面便點「投屏」按鈕）。
  \item 投屏成功後，試嚇喺手機上做其他操作，睇嚇大屏係咪同步顯示。
\end{enumerate}

\end{tcbpractice}

\begin{tcbsummary}{本章小結}
\begin{itemize}
  \item 無線投屏三陣營：AirPlay（蘋果）、Miracast（安卓）、DLNA（跨平臺推媒體）
  \item 有線投屏更穩定零延遲，會議室場景必備轉HDMI線
  \item Samsung DeX/華為桌面模式等手機變身輕辦公電腦
  \item 車載互聯：CarPlay/Android Auto/HiCar/CarLife四大方案
  \item 會議室投屏鐵律：提前測試、準備Plan B、關閉通知預覽
\end{itemize}

\end{tcbsummary}
